\chapter{相关理论与方法基础}

\section{标准数据分析流程}

标准数据分析流程通常包括问题定义、数据获取、数据质量检查、特征工程、探索性分析、模型训练、模型验证、结果解释和复现归档。本文按照该流程组织实验：先确认数据集满足课程对年份、规模、格式和任务类型的要求，再进行预处理与可视化，最后分别完成分类、回归和聚类建模。这样做的目的，是保证模型结果不是孤立指标，而是可以追溯到数据来源、变量选择和实验设置。

\section{分类任务与评价指标}

分类任务的目标是学习输入特征 $X$ 与离散标签 $y$ 之间的映射。本文分类任务预测 \texttt{high\_risk\_flag}，属于二分类筛查问题。除 Accuracy 外，本文更关注 Precision、Recall、F1、PR-AUC 和混淆矩阵，因为高风险筛查中漏判会降低模型的参考价值。

若 $TP$、$FP$、$TN$、$FN$ 分别表示真正例、假正例、真负例和假负例，则主要指标定义为：
\begin{align}
\mathrm{Precision} &= \frac{TP}{TP+FP}, \\
\mathrm{Recall} &= \frac{TP}{TP+FN}, \\
\mathrm{F1} &= \frac{2 \times \mathrm{Precision} \times \mathrm{Recall}}{\mathrm{Precision}+\mathrm{Recall}}.
\end{align}

Precision 反映被模型标为高风险的样本中有多少确为高风险，Recall 反映真实高风险样本中有多少被识别出来。F1 是 Precision 与 Recall 的调和平均，用于在二者之间取得综合评价。

\section{回归任务与评价指标}

回归任务用于预测连续型目标变量。本文以 \texttt{digital\_dependence\_score} 为主回归目标，同时保留 \texttt{productivity\_score} 作为弱预测对照。评价指标包括 $R^2$、MSE、RMSE 和 MAE：
\begin{align}
\mathrm{MSE} &= \frac{1}{n}\sum_{i=1}^{n}(y_i-\hat{y}_i)^2, \\
\mathrm{MAE} &= \frac{1}{n}\sum_{i=1}^{n}|y_i-\hat{y}_i|, \\
R^2 &= 1-\frac{\sum_{i=1}^{n}(y_i-\hat{y}_i)^2}{\sum_{i=1}^{n}(y_i-\bar{y})^2}.
\end{align}

MSE 与 RMSE 对较大误差更敏感，MAE 更直观地表示平均绝对偏差。$R^2$ 衡量模型相对均值预测的解释能力，若 $R^2$ 接近 0 或为负，说明当前特征和模型难以稳定预测该目标。

\section{聚类分析与 PCA 降维}

聚类分析不使用监督标签，而是依据样本在特征空间中的相似性进行分组。本文比较 KMeans、AgglomerativeClustering 和 GaussianMixture，并使用肘部法则、Silhouette、Calinski-Harabasz 和 Davies-Bouldin 等指标辅助判断。

Silhouette 衡量样本与本簇和邻近簇之间的相对距离。对样本 $i$，设 $a(i)$ 为其到同簇样本的平均距离，$b(i)$ 为其到最近其他簇样本的平均距离，则：
\begin{equation}
s(i)=\frac{b(i)-a(i)}{\max\{a(i),b(i)\}}.
\end{equation}

Silhouette 越高通常表示簇内更紧凑、簇间更分离。本文的聚类 Silhouette=0.1860，说明分群边界偏弱，因此聚类结果只作为探索性用户画像。

PCA 用于将高维标准化特征投影到少数主成分上，便于可视化和辅助理解。本文的 PCA 不作为分类或回归模型输入，也不用于因果解释。

\section{目标泄漏控制与模型验证原则}

目标泄漏指模型训练时使用了与目标变量高度重合或由目标派生的信息，导致测试指标虚高但实际不可用。本文在分类预测 \texttt{high\_risk\_flag} 时，明确排除 \texttt{anxiety\_score}、\texttt{depression\_score}、\texttt{stress\_level}、\texttt{happiness\_score}、\texttt{focus\_score}、\texttt{productivity\_score} 和 \texttt{digital\_dependence\_score} 等 outcome 字段。回归任务也排除 \texttt{high\_risk\_flag}、心理状态结果变量和另一个 outcome 目标变量。聚类任务只使用数字行为与生活习惯数值特征，将背景类别和结果变量仅用于聚类后的画像解释。

模型验证方面，分类任务使用 StratifiedKFold、超参数搜索和独立测试集评估；阈值选择在验证集上完成，最终测试集只评估一次。回归任务使用交叉验证选择超参数，并在测试集报告最终指标。聚类任务不使用监督标签训练，而是用多种内部指标和可视化结果共同判断分群是否稳定。
